\section{Laman graphs}
\label{sec:laman}

This section specializes \cref{sec:sparsity} to $(k,\ell) = (2,3)$,
the case of Laman graphs. Pollaczek-Geiringer~\cite{pollaczekGeiringer1927}
(pp.~60 and~64), and later Laman~\cite{laman1970}, proved that the Laman
graphs are exactly the generically minimally rigid graphs in the plane
(see Jord\'an~\cite{jordan2016}, Theorem~2.2.2). Equivalently, a graph
on at least two vertices is generically rigid in the plane if and only
if it has a spanning Laman subgraph. We prove this form in
\cref{sec:laman-theorem}.

The direction from Laman graphs to rigidity is proved by induction on
the number of vertices. Its base case is a Laman graph on two vertices,
and its inductive step deletes a vertex of degree two or three by
reversing a Henneberg move (\cref{sec:henneberg}). This section proves
the facts about Laman graphs that the induction rests on: the Laman
property is invariant under isomorphism, the only Laman graph on two
vertices is the complete graph $K_2$, and in a Laman graph on at least
three vertices every vertex has degree at least two.

\subsection{Definition and isomorphism invariance}

\begin{definition}[Laman graph]
  \label{def:isLaman}
  \lean{SimpleGraph.IsLaman}
  \leanok
  \uses{def:isTight}
  Let $V$ be finite. A graph $G$ on $V$ is \emph{Laman} if it is
  $(2,3)$-tight, that is, if it is $(2,3)$-sparse and
  \[
    |E(G)| + 3 \;=\; 2\,|V|.
  \]
\end{definition}

\begin{lemma}[Laman is preserved by isomorphism]
  \label{lem:isLaman-iso}
  \lean{SimpleGraph.IsLaman.iso}
  \leanok
  \uses{def:isLaman, lem:isTight-iso}
  Let $\varphi$ be an isomorphism from $G$ to a graph $H$. If $G$ is
  Laman, then so is $H$.
\end{lemma}
\begin{proof}
  \leanok
  Specialize \cref{lem:isTight-iso} to $(k, \ell) = (2, 3)$.
\end{proof}

\subsection{The base case: \texorpdfstring{$K_2$}{K2}}

The first result below is the starting point of the construction of
Laman graphs by Henneberg moves, as in the example of
\cref{thm:top-fin-four-minus-edge-isLaman}. The second is the base case
of the induction in the proof of Laman's theorem
(\cref{thm:isLaman-isGenericallyRigid-two}).

\begin{theorem}[$K_2$ is Laman]
  \label{thm:top-fin-two-isLaman}
  \lean{SimpleGraph.top_fin_two_isLaman}
  \leanok
  \uses{def:isLaman, lem:edgesIn-corners}
  The complete graph $K_2$ on two vertices is Laman.
\end{theorem}
\begin{proof}
  \leanok
  The graph $K_2$ has one edge, and $1 + 3 = 2 \cdot 2$. For sparsity,
  the only set $s$ of vertices with $3 \le 2\,|s|$ is the whole vertex
  set. It spans the one edge (\cref{lem:edgesIn-corners}), and
  $1 + 3 \le 2 \cdot 2$.
\end{proof}

\begin{lemma}[Laman graphs on two vertices are $K_2$]
  \label{lem:isLaman-eq-top-of-card-eq-two}
  \lean{SimpleGraph.IsLaman.eq_top_of_card_eq_two}
  \leanok
  \uses{def:isLaman}
  Let $G$ be a Laman graph on a set $V$ of two vertices. Then $G$ is the
  complete graph on $V$.
\end{lemma}
\begin{proof}
  \leanok
  Since $|E(G)| + 3 = 2 \cdot 2$, the graph $G$ has one edge. The
  complete graph on $V$ also has exactly one edge, and it contains
  every edge of $G$. So the two edge sets are equal.
\end{proof}

\subsection{Degree bounds}
\label{sec:laman-degrees}

Both lemmas below concern a Laman graph on at least three vertices.
The reverse decomposition uses the first to exclude a vertex of degree
one (\cref{thm:isLaman-exists-typeI-or-typeII-reverse}). Lee and
Streinu~\cite[Lemma~4]{leeStreinu2008} give the same argument for
$(k,\ell)$-tight graphs. The second lemma is the Laman form of
\cref{lem:isSparse-exists-degree-le-three}.

\begin{lemma}[Minimum degree]
  \label{lem:isLaman-two-le-degree}
  \lean{SimpleGraph.IsLaman.two_le_degree}
  \leanok
  \uses{def:isLaman, lem:ncard-edgesIn-compl-singleton}
  Let $V$ be finite with $|V| \ge 3$, and let $G$ be a Laman graph on
  $V$. Then every vertex of $G$ has degree at least $2$.
\end{lemma}
\begin{proof}
  \leanok
  Let $v \in V$. The set $\{v\}^c$ has $|V| - 1 \ge 2$ vertices, so
  sparsity applies to it:
  $|\edgesIn{\{v\}^c}| + 3 \le 2(|V| - 1)$. By
  \cref{lem:ncard-edgesIn-compl-singleton},
  $|\edgesIn{\{v\}^c}| = |E(G)| - \deg_G(v)$. Since
  $|E(G)| + 3 = 2|V|$, this gives $2|V| - \deg_G(v) \le 2|V| - 2$, that
  is, $\deg_G(v) \ge 2$.
\end{proof}

\begin{lemma}[A vertex of degree 2 or 3]
  \label{lem:isLaman-exists-two-le-degree-le-three}
  \lean{SimpleGraph.IsLaman.exists_two_le_degree_le_three}
  \leanok
  \uses{lem:isLaman-two-le-degree, lem:isSparse-exists-degree-le-three}
  Let $V$ be finite with $|V| \ge 3$, and let $G$ be a Laman graph on
  $V$. Then some vertex of $G$ has degree $2$ or $3$.
\end{lemma}
\begin{proof}
  \leanok
  Since $G$ is $(2,3)$-sparse, it has a vertex of degree at most $3$
  (\cref{lem:isSparse-exists-degree-le-three}). That vertex has degree
  at least $2$ (\cref{lem:isLaman-two-le-degree}).
\end{proof}
